% =============================================================================
% Title + problem + motivation (slides 1–4, ~3.5 min)
% =============================================================================

% --- Slide 1: Title ---
\begin{frame}[plain,noframenumbering]
\vspace*{0.4em}
\begin{columns}[T,onlytextwidth]
\begin{column}{0.68\textwidth}
\raggedright
{\color{uaibluedark}\fontsize{18}{22}\selectfont\bfseries
Geometric Filtering of LLM-Generated Samples for Few-Shot Text Classification\par}

\vspace{0.6em}
{\color{uaiorange}\rule{0.55\textwidth}{1pt}}
\vspace{0.8em}

{\large\bfseries\color{uaibluedark} Benjamín Schindler \quad Gonzalo A. Ruz}\\[4pt]
{\footnotesize\color{uaigray}
Universidad Adolfo Ibáñez, Santiago, Chile\\
Millennium Nucleus for Social Data Science (SODAS)\\
Millennium Nucleus in Data Science for Plant Resilience (PhytoLearning)\par}

\vspace{1.2em}
{\small\bfseries\color{uaibluedark} IEEE LACCI 2026}\\[2pt]
{\footnotesize\color{uaigray} Universidad del Pacífico \textperiodcentered{} Lima, Peru \textperiodcentered{} November 2026}
\end{column}
\begin{column}{0.30\textwidth}
\centering
\vspace{0.4em}
\includegraphics[height=1.2cm]{logos/fic.png}

\vspace{2.6em}
{\scriptsize\color{uaigray}
benjamin.schindler@alumnos.uai.cl\\
gonzalo.ruz@uai.cl\par}
\end{column}
\end{columns}
\end{frame}

% --- Slide 2: The problem (shared with the simple deck) ---
% =============================================================================
% Intro slide shared by both decks (slide 2): why few labels are a problem and
% why generated data must be reliable. Input once per document.
% =============================================================================

% Toy 2-D sketch: class A (left) vs. class B (right); the true boundary is x = 0.
\newcommand{\fewshotpanel}{%
  \draw[uaigray!50, dashed] (0,-1.15) -- (0,1.15);
  \foreach \pt in {(-1.0,0.6),(-0.7,-0.5),(-1.3,-0.1)} {\node[real] at \pt {};}
  \foreach \pt in {(0.9,0.4),(1.2,-0.6),(0.6,-0.2)} {\node[realB] at \pt {};}
}
\begin{frame}{The problem: learning from very few examples}
\small
Many real tasks have only \textbf{10--50 labeled examples per class}: labeling is slow and costly.

\vspace{0.6em}
\begin{columns}[T,onlytextwidth]
\begin{column}{0.32\textwidth}
\centering
{\bfseries\color{uaibluedark} 1.\ Few examples}\\[4pt]
\begin{tikzpicture}
  \fill[uaigraylight!60] (-1.6,-1.15) rectangle (1.6,1.15);
  \fewshotpanel
  \draw[uaiorange, thick] (-0.5,-1.15) -- (0.4,1.15);
  \draw[uaiorange, thick] (0.4,-1.15) -- (-0.55,1.15);
  \draw[uaiorange, thick] (0.15,-1.15) -- (0.2,1.15);
  \node[font=\large\bfseries, color=uaiorange, fill=uaigraylight, circle, inner sep=0.5pt] at (0.0,0.0) {?};
\end{tikzpicture}\\[4pt]
Many boundaries fit the data.\\
\textbf{Which one is right?}
\end{column}
\begin{column}{0.32\textwidth}
\centering
{\bfseries\color{uaibluedark} 2.\ Create more from them}\\[4pt]
\begin{tikzpicture}
  \fill[uaigraylight!60] (-1.6,-1.15) rectangle (1.6,1.15);
  \fewshotpanel
  \foreach \pt in {(-0.4,0.9),(-0.3,-0.2),(-1.4,0.8),(-0.9,-0.95),(-0.25,0.4)} {\node[genA] at \pt {};}
  \foreach \pt in {(0.3,0.7),(0.35,-0.8),(1.4,0.2),(0.8,0.95),(0.25,0.1)} {\node[genB] at \pt {};}
  \draw[boundary] (0.02,-1.15) -- (-0.02,1.15);
\end{tikzpicture}\\[4pt]
Reliable new examples\\
\textbf{pin the boundary down.}
\end{column}
\begin{column}{0.32\textwidth}
\centering
{\bfseries\color{uaibluedark} 3.\ But only if reliable}\\[4pt]
\begin{tikzpicture}
  \fill[uaigraylight!60] (-1.6,-1.15) rectangle (1.6,1.15);
  \fewshotpanel
  \foreach \pt in {(-0.4,0.9),(-0.3,-0.2),(-1.4,0.8)} {\node[genA] at \pt {};}
  \foreach \pt in {(1.4,0.2),(0.8,0.95)} {\node[genB] at \pt {};}
  \foreach \pt in {(0.45,-0.75),(0.55,0.1),(0.35,0.65)} {\node[genBad] at \pt {};}
  \draw[boundary=methodWorse] (1.05,-1.15) -- (0.45,1.15);
\end{tikzpicture}\\[4pt]
Wrong examples\\
\textbf{teach the wrong boundary.}
\end{column}
\end{columns}

\vspace{0.5em}
\centering
\begin{tikzpicture}
  \node[real] at (0,0) {};    \node[legendtxt] at (0.12,0) {real A};
  \node[realB] at (1.4,0) {}; \node[legendtxt] at (1.52,0) {real B};
  \node[genA] at (2.8,0) {};  \node[legendtxt] at (2.92,0) {new A};
  \node[genB] at (4.15,0) {}; \node[legendtxt] at (4.27,0) {new B};
  \node[genBad] at (5.5,0) {}; \node[legendtxt] at (5.62,0) {new ``A'' that is really B};
  \draw[uaigray!60, dashed] (9.0,0) -- (9.4,0); \node[legendtxt] at (9.45,0) {true boundary};
\end{tikzpicture}

\keymessage{We need more training data, but only data we can trust.}
\end{frame}


% --- Slide 3: Two ways to create data ---
\begin{frame}{Two ways to create more training data}
\small
\begin{columns}[T]
\begin{column}{0.5\textwidth}
\centering
\setbeamercolor{block title}{fg=white,bg=methodSMOTE}
\begin{block}{Classical oversampling (SMOTE)}
\scriptsize
$x_{\text{new}} = x_i + \lambda\,(x_j - x_i)$ in embedding space\\[3pt]
\good\ Geometric control\\
\bad\ Produces \alert{vectors, not text}; only fills known space
\end{block}
\begin{tikzpicture}[scale=0.85]
  \embeddingclasses
  % interpolation segments between real neighbors of class A
  \foreach \a/\b in {pa1/pa2, pa2/pa3, pa4/pa3, pa5/pa4, pa1/pa5} {\draw[uaiorange!70, dashed] (\a) -- (\b);}
  \foreach \a/\b/\t in {pa1/pa2/0.45, pa2/pa3/0.6, pa4/pa3/0.4, pa5/pa4/0.55, pa1/pa5/0.5} {\node[synth] at ($(\a)!\t!(\b)$) {};}
  \node[font=\scriptsize, color=uaigray] at (1.25,-1.3) {stays on segments between real neighbors};
\end{tikzpicture}
\end{column}
\begin{column}{0.5\textwidth}
\centering
\setbeamercolor{block title}{fg=white,bg=methodOurs}
\begin{block}{LLM generation (GPT, Gemini, Claude\dots)}
\scriptsize
Prompt with real examples, ask for new ones\\[3pt]
\good\ Fluent, diverse, valid text\\
\bad\ \alert{No geometric awareness}: noise, class overlap
\end{block}
\begin{tikzpicture}[scale=0.85]
  \embeddingclasses
  % LLM samples meant for class A: some inside, some peripheral, some cross-class
  \foreach \pt in {(-0.15,0.15),(0.25,-0.35),(-0.45,0.65)} {\node[synth] at \pt {};}
  \foreach \pt in {(-1.35,0.75),(-1.2,-0.85),(0.85,1.05)} {\node[synth] at \pt {};}
  \foreach \pt in {(2.45,0.05),(2.75,0.55)} {\node[synth] at \pt {};}
  \node[font=\tiny\itshape, color=uaiorange, anchor=south] at (2.6,0.85) {cross-class};
  \node[font=\tiny\itshape, color=uaiorange, anchor=east] at (-1.42,0.75) {peripheral};
  \node[font=\scriptsize, color=uaigray] at (1.25,-1.3) {open-ended: no notion of the class region};
\end{tikzpicture}
\end{column}
\end{columns}

\centering
\begin{tikzpicture}
  \node[real] (l1) at (0,0) {};            \node[legendtxt] at (0.12,0) {real, class A};
  \node[realB] (l2) at (2.3,0) {};         \node[legendtxt] at (2.42,0) {real, class B};
  \node[synth] (l3) at (4.6,0) {};         \node[legendtxt] at (4.72,0) {synthetic sample for class A};
\end{tikzpicture}

\keymessage{Can we get LLM-quality text \emph{and} SMOTE-like geometric control?}
\end{frame}

% --- Slide 4: The idea ---
\begin{frame}{The idea: filter LLM samples in embedding space}
\begin{columns}[T]
\begin{column}{0.54\textwidth}
\small
LLM samples have \textbf{heterogeneous quality}: some land inside their class region,
others in peripheral or cross-class zones. Adding all of them can \alert{hurt} the classifier.

\vspace{0.5em}
\textbf{Post-generation geometric filtering:}
\begin{enumerate}\setlength\itemsep{1pt}\footnotesize
  \item Generate a \textbf{3$\times$ surplus} of candidates with an LLM.
  \item Embed them with the same sentence encoder.
  \item Score each candidate by its \textbf{distance to real examples} of its class.
  \item Keep the best; \textbf{weight} them by their score.
\end{enumerate}

\vspace{0.4em}
\begin{insight}
\footnotesize
\textbf{Research question:} can geometric filtering make LLM augmentation
reliably better than SMOTE in few-shot settings?
\end{insight}
\end{column}
\begin{column}{0.46\textwidth}
\centering
\vspace{0.6em}
\begin{tikzpicture}[scale=1.15, every node/.append style={transform shape=false}]
  \embeddingclasses
  \draw[uaibluedark, dashed, thick] (0,0) circle (1.45);
  % candidates generated for class A
  \foreach \pt in {(-0.15,0.1),(0.35,0.2),(-0.2,0.85),(0.7,0.65),(-0.95,0.0),(0.55,-0.75),(-0.85,-0.75),(1.15,-0.3)} {\node[kept] at \pt {};}
  \foreach \pt in {(-1.75,0.75),(-1.45,-1.25),(0.3,1.75),(2.4,0.05),(2.75,0.65),(1.75,-1.15)} {\node[rejected] at \pt {};}
  % distance of one rejected candidate to its nearest real example
  \draw[<->, >=Stealth, uaigray, thin, shorten <=2pt, shorten >=3pt] (pa1) -- (-1.75,0.75);
  \node[font=\scriptsize, color=uaigray] at (-1.0,0.8) {$d(x)$};
  \node[font=\scriptsize, color=uaibluedark, anchor=south west] at (0.75,1.2) {keep top-$N$};
  \node[font=\tiny\itshape, color=uaiorange, anchor=north] at (2.6,-0.62) {cross-class};
\end{tikzpicture}

\vspace{0.6em}
\begin{tikzpicture}
  \node[real]     at (0,0.45) {};  \node[legendtxt] at (0.12,0.45) {real example (class A / B)};
  \node[realB]    at (-0.16,0.45) {};
  \node[kept]     at (0,0)    {};  \node[legendtxt] at (0.12,0)    {LLM candidate, kept};
  \node[rejected] at (0,-0.45) {}; \node[legendtxt] at (0.12,-0.45) {LLM candidate, rejected};
\end{tikzpicture}
\end{column}
\end{columns}
\end{frame}
